\section{Определение интеграла от функции комплексного переменного по
ориентированному контуру. Свойства интеграла от функции комплексного
переменного.}\label{ux43eux43fux440ux435ux434ux435ux43bux435ux43dux438ux435-ux438ux43dux442ux435ux433ux440ux430ux43bux430-ux43eux442-ux444ux443ux43dux43aux446ux438ux438-ux43aux43eux43cux43fux43bux435ux43aux441ux43dux43eux433ux43e-ux43fux435ux440ux435ux43cux435ux43dux43dux43eux433ux43e-ux43fux43e-ux43eux440ux438ux435ux43dux442ux438ux440ux43eux432ux430ux43dux43dux43eux43cux443-ux43aux43eux43dux442ux443ux440ux443.-ux441ux432ux43eux439ux441ux442ux432ux430-ux438ux43dux442ux435ux433ux440ux430ux43bux430-ux43eux442-ux444ux443ux43dux43aux446ux438ux438-ux43aux43eux43cux43fux43bux435ux43aux441ux43dux43eux433ux43e-ux43fux435ux440ux435ux43cux435ux43dux43dux43eux433ux43e.}

\textbf{Кривые:}

\begin{itemize}
\tightlist
\item
  \(\gamma:\begin{cases}x=x(t) \\ y=y(t)\end{cases}, \ t\in[\alpha;\beta], \ x(t),y(t)\in C^{1}[\alpha;\beta]\)
  - уравнение гладкой кривой;
\item
  \(\gamma:x(\alpha)=x(\beta), \ y(\alpha)=y(\beta)\) - замкнутая кривая
\item
  \(\gamma:\) Кривая Жордана: из
  \(t_{1}\neq t_{2}\in[\alpha;\beta)\implies(x(t_{1}),y(t_{1}))\neq(x(t_{2}),y(t_{2}))\)
  - простая кривая
\item
  \(\gamma:\) ориентированная кривая:

  \begin{itemize}
  \tightlist
  \item
    указание на начало и конец \(\gamma\) (незамкнутая кривая)
  \item
    направление обхода границы области \(G: \ \gamma=\partial G\)
  \item
    по возрастанию параметр \(t\)
  \end{itemize}
\end{itemize}

\textbf{Определение:} Интеграл функции \(f(z)\) по \(\gamma:\)

\begin{enumerate}
\def\labelenumi{\arabic{enumi}.}
\tightlist
\item
  \(p\) - разбиение \([\alpha;\beta]\), \(t=\{ t_{i} \}_{i=1}^{n}\) -
  точки разбиения \(p\), \(\Delta t_{i}=t_{i}-t_{i-1}\) - отрезок,
  \(L(p)=\max\limits_{i\in \{ 1,\dots,n \}}\Delta t_{i}\) - мелкость
  разбиения, набор \(\xi=\{ \xi_{i} \}_{i=1}^{n}\) - выборка из
  разбиения \(p\), если
  \(\forall i=\overline{1,n}\quad \xi_{i}\in H_{i}:=[t_{i-1},t_{i}]\)
\item
  \(\Delta z_{i}=z(t_{i})- z(t_{i-1}), \ z=z(t)\) - уравнение
  \(\gamma, \ |\Delta z_{i}|=|\Delta l_{i}|\)
\item
  \(\sigma(f,p,\xi)=\sum\limits_{i=1}^{n}f(z(\xi_{i}))\cdot\Delta z_{i}\)
  - интегральная сумма
\end{enumerate}

\textbf{Определение:} Интегралом функции \(f(z)\) по \(\gamma\),
обозначение \(I=\int\limits_{\gamma}f(z)\,dz\), называют число
\(I\in \mathbb{C}: \ \lim_{ L(p) \to 0}\sigma(f,p,\xi)=I\): \[
\forall\varepsilon>0 \ \exists\delta>0 \forall p \in P_{\delta} \ \forall \xi \in V_{p}\implies |\sigma(f,p,\xi)-I|<\varepsilon
\]

Функция называется интегрируемой на \(\gamma\), если
\(\exists \int\limits_{\gamma}f(z)\,dz\).

\textbf{Свойства:}

\begin{enumerate}
\def\labelenumi{\arabic{enumi}.}
\tightlist
\item
  Линейность \(\forall f(z), g(z)\) - интегрируемые функции на
  \(\gamma\), \(\forall\lambda,\mu \in \mathbb{C}\). Тогда \[
  \int_{\gamma}(\lambda\cdot f(z)+\mu\cdot g(z))\,dz=\lambda \int_{\gamma}f(z)\,dz+\mu \int_{\gamma}g(z)\,dz
  \]
\item
  \(\gamma^{+}\) и \(\gamma^{-}\) - кривые с разной ориентацией,
  \(f(z)\) - интегрируемая функция на \(\gamma^{+}, \gamma^{-}\): \[
  \int_{\gamma^{+}}f(z)\,dz=-\int_{\gamma^{-}}f(z)\,dz
  \]
\item
  Аддитивность интеграла по множеству:
  \(\gamma=\gamma_{1}\cup\gamma_{2}, \ f(z)\) интегрируема на
  \(\gamma_{1}\) и \(\gamma_{2}\). Тогда \[
  \int_{\gamma_{1}\cup\gamma_{2}}f(z)\,dz=\int_{\gamma_{1}}f(z)\,dz+\int_{\gamma_{2}}f(z)\,dz
  \]
\item
  Оценка интеграла: \(\gamma:z=z(t), \ t\in[\alpha;\beta]\) \[
  \left\lvert \int f(z)\,dz   \right\rvert \leq \int_{\gamma} |f(z)|\,d\ell \overset{\mathrm{def}}{=}\int_{\gamma}|f(z)|\,|dz|,
  \] где
  \(d \ell=\sqrt{ x'^{2}(t)+y'^{2}(t) }\,dt=|z'(t)|\,dt=|z'(t)\,dt|=|dz|\)
\item
  Если ряд \(f(z)=\sum\limits_{n=1}^{\infty}f_{n}(z)\), составленный из
  непрерывных на кривой \(\gamma\) функций \(f_{n}(z)\)
  \((n=1,2,\dots)\), сходится равномерно на \(\gamma\), то его можно
  почленно интегрировать, т.е. \[
  \int_{\gamma}f(z)\,dz=\sum_{n=1}^{\infty}\int_{\gamma}f_{n}(z)\,dz.
  \]
\end{enumerate}
